\documentclass[11pt,a4paper]{article}
\usepackage[utf8]{inputenc}
\usepackage[T1]{fontenc}
\usepackage{estilo_capstone}

\titulo{Plan de pruebas — Evento Cultural}
\subtitulo{Plan de pruebas y evidencias (artefacto 7 del marco ágil, anexo metodológico)}
\autor{Roberto Riffo · José Miguel Moncada (QA) · David Herrera}
\fecha{25 de septiembre de 2026}

\newcolumntype{L}[1]{>{\raggedright\arraybackslash}p{#1}}
\definecolor{rojo}{HTML}{B03A3A}   % misma definición del documento técnico

\begin{document}
\maketitlecapstone

\section{Propósito y alcance}

Este plan organiza la verificación del proyecto según la exigencia del anexo
metodológico (pruebas unitarias, de integración, de rendimiento y de seguridad)
y la relaciona con los requisitos no funcionales priorizados en Fase 1
(RNF-01 a RNF-12) y con la verificación ya ejecutada. Declara el estado real:
lo ejecutado con su resultado, y lo pendiente sin asumirlo como hecho.

\section{Estado actual de las pruebas}

\begin{table}[h!]
\centering\small
\begin{tabular}{@{}L{5.2cm}L{4.6cm}L{5.0cm}@{}}
\toprule
\textbf{Suite} & \textbf{Alcance} & \textbf{Comando} \\
\midrule
Pruebas unitarias del scraper (14 módulos) & Conectores, normalización, filtro cultural, deduplicación, checkpoints y exportación & \code{python -m unittest discover -s tests -v} \\
Pruebas del frontend (Vitest) & Servicios y lógica de componentes & \code{npm test} \\
\bottomrule
\end{tabular}
\caption{Suites de pruebas existentes y su ejecución.}
\end{table}

\section{Cobertura por requisito no funcional}

\begin{longtable}{@{}L{1.6cm}L{4.4cm}L{3.6cm}L{5.2cm}@{}}
\toprule
\textbf{RNF} & \textbf{Requisito} & \textbf{Tipo de prueba} & \textbf{Estado} \\
\midrule
\endhead
RNF-01 & Interfaz clara e intuitiva & Usabilidad & Pendiente: evaluación con usuarios (ver §5) \\
RNF-02 & Diseño responsive / multiplataforma & Funcional manual & \textbf{Ejecutada}: desbordes revisados en 320--1440 px; verificación manual en navegador y emulador de dispositivos \\
RNF-03 & Tiempos de carga aceptables & Rendimiento & Parcial: skeletons y estados de carga implementados; medición formal pendiente \\
RNF-04 & Protección de datos de usuarios y organizadores & Seguridad & Pendiente en frontend (no hay datos personales: estado mock); RLS de Supabase documentada, validación de políticas pendiente \\
RNF-04 & Control de acceso por rol & Seguridad & Pendiente: autenticación real no conectada \\
RNF-06 & Disponibilidad del servicio & Disponibilidad & Parcial: workflow de GitHub Actions operativo con ejecución manual y cron diario activo desde el 24-09 (verificado en GitHub Actions: corridas \#9--\#11 del 25 al 27-09); pendiente una ejecución automática completa tras corregir la prueba dependiente del tiempo (§5.1) \\
RNF-07 & Soporte de crecimiento de datos & Rendimiento & Parcial: paginación/búsqueda indexada no medida con volumen real \\
RNF-08 & Consistencia y normalización & Funcional automatizada & \textbf{Ejecutada}: pruebas del pipeline (normalización y deduplicación) en la suite del scraper \\
RNF-09 & Código mantenible y documentado & Estático & \code{npm run lint} sin errores; lint de Python por revisión manual \\
RNF-10 & Contenido en español y contexto chileno & Funcional manual & \textbf{Ejecutada}: interfaz en es-CL; catálogo con comunas y regiones de Chile \\
RNF-11 & Despliegue reproducible & Integración & En curso: workflow del scraper desplegado y en testeo del despliegue (evidencia de GitHub Actions); alojamiento web en Vercel y conversión a app con Capacitor en preparación \\
RNF-11/12 & Accesibilidad web & Accesibilidad & Parcial: contraste AA/AAA calculado, teclado y foco verificados en navegador; prueba con lectores de pantalla pendiente \\
\bottomrule
\end{longtable}

\section{Pruebas pendientes según el anexo}

\begin{table}[h!]
\centering\small
\begin{tabular}{@{}L{3.4cm}L{11.4cm}@{}}
\toprule
\textbf{Tipo (anexo)} & \textbf{Plan para cubrirla} \\
\midrule
Integración & End-to-end scraper$\rightarrow$Supabase$\rightarrow$frontend: ejecutar el pipeline con sincronización a Supabase y validar el catálogo resultante; y conexión real del frontend a Supabase (hoy mock), verificando el contrato de \code{recomendacion.model.ts} \\
Rendimiento & Medir tiempos de carga del frontend con el catálogo sincronizado y la consulta a Supabase en volumen; repetir \code{cargarEventos} con datos reales \\
Seguridad & Validar las políticas RLS de Supabase con roles reales (anónimo, autenticado) y revisar el manejo de secretos del workflow \\
Usabilidad (RNF-01) & Evaluación con usuarios objetivo; puede integrarse con la validación de historias de usuario HU-01--15 \\
\bottomrule
\end{tabular}
\caption{Pruebas pendientes y cómo se cubrirán.}
\end{table}

\subsection{Planificación de las pruebas pendientes}

La planificación se organiza en tres bloques en cascada, porque cada uno habilita al
siguiente; el orden respeta la situación real del proyecto (datos de ejemplo en el
frontend, integración con Supabase pendiente):

\begin{enumerate}[leftmargin=1.6em]
    \item \textbf{Bloque 1 --- Integración de datos (habilita a los demás):}
    sincronización scraper$\rightarrow$Supabase con validación del catálogo resultante
    (RNF-06/08); conexión real del frontend a Supabase sustituyendo
    \code{data/mock-eventos.ts}, verificando el contrato de \code{recomendacion.model.ts}.
    \emph{Criterio de salida:} el catálogo consultado por el frontend proviene de
    Supabase y las pruebas de la suite del scraper pasan contra el entorno real.
    \item \textbf{Bloque 2 --- Medición sobre datos reales:} tiempos de carga del
    frontend (RNF-03), consulta a Supabase con volumen real (RNF-07) y validación de
    políticas RLS con roles reales (RNF-04/05). \emph{Criterio de salida:} métricas
    registradas en la §5 con fecha y volumen de eventos.
    \item \textbf{Bloque 3 --- Validación con usuarios y accesibilidad:} evaluación de
    usabilidad (RNF-01) integrada con la revisión de HU-01--15, y prueba con lectores
    de pantalla (RNF-12). \emph{Criterio de salida:} acta de resultados por historia y
    registro de hallazgos de accesibilidad.
\end{enumerate}

Los bloques 1 y 2 dependen del despliegue/estado del catálogo real; el bloque 3 puede
planificarse en paralelo al 2. Cada ejecución se agrega a la §5 sin reescribir el
histórico, y el responsable de cada bloque queda \textcolor{capstoneprimario}{\textbf{por
asignar por el equipo}} (roles de referencia: BD/QA para bloques 1 y 2, Frontend/UX para
el 3).

\section{Registro de resultados}

Los resultados de las pruebas ya ejecutadas están documentados en la
documentación técnica (\code{Documento\_Tecnico\_Capstone}, sección de pruebas)
y en el control de cambios del frontend. Este plan se actualizará con cada
ronda de pruebas; los resultados nuevos se agregan sin reescribir el histórico.

\subsection{Incidencia de integración continua (27-09-2026)}

Fallo registrado en GitHub Actions durante la ejecución automática de la suite del
scraper (corrida \#11, \code{scrape-events.yml}, disparo programado del cron diario,
27-09-2026 09:12 UTC, commit \code{63e310c}; verificado contra la API pública de
GitHub Actions). Se consigna aquí como registro detallado del incidente (no existe un
artefacto separado de incidencias en el repositorio).

\begin{table}[h!]
\centering\small
\begin{tabular}{@{}L{3.4cm}L{11.4cm}@{}}
\toprule
\textbf{Campo} & \textbf{Detalle} \\
\midrule
Fecha & 27-09-2026 \\
Componente & Scraper de eventos culturales \\
Entorno & GitHub Actions (integración continua), corrida \#11 de \code{scrape-events.yml} \\
Tipo & Fallo de prueba automatizada dependiente del tiempo \\
Severidad & Baja \\
Impacto & Ninguno en producción; bloqueo del workflow antes de ejecutar la extracción (pasos de extracción y sincronización saltados) \\
Prueba afectada & \code{test\_ceina\_cards\_are\_mapped\_to\_canonical\_event\_fields} \\
Resultado & 138 pruebas ejecutadas: 137 correctas, 1 fallida (\code{AssertionError: 0 != 1}) \\
Causa raíz & La prueba simulaba un evento de CEINA entre el 24 y el 26-09-2026 y congelaba la fecha del parser en el 21-09 (\code{chile\_now}); la comprobación de vigencia \code{is\_upcoming()} de \code{sanitization.py} no estaba bajo ese reloj simulado y usó la fecha real de ejecución (27-09), descartando el evento como vencido. Verificado en el código fuente del repositorio: \code{tests/test\_event\_sources.py} y \code{eventos/sanitization.py} \\
Aclaración & No se evidenció caída de CEINA ni fallo de Supabase, NVIDIA o del proceso real de extracción; los mensajes \code{ERROR:root} del registro eran excepciones simuladas por pruebas que finalizaron correctamente \\
Acción correctiva & Congelar en la prueba tanto la fecha utilizada por el parser como la utilizada por \code{is\_upcoming()} (reloj controlado); recomendación extendida a toda prueba que evalúe vigencia \\
Estado & \textcolor{rojo}{\textbf{Pendiente de corrección}}: no se marca resuelto sin evidencia de una ejecución posterior exitosa (la corrida del 26-09, \#10, sí pasó; la falla específica surgió al vencer las fechas simuladas) \\
Criterios de cierre & (1) la prueba específica termina correctamente; (2) las 138 pruebas terminan correctamente; (3) GitHub Actions finaliza con código de salida 0 \\
Aprendizaje & Las pruebas que evalúan vigencia de eventos deben utilizar un reloj controlado y no depender de la fecha real del sistema \\
\bottomrule
\end{tabular}
\caption{Incidencia de CI: prueba dependiente del tiempo en la suite del scraper.}
\label{tab:incidencia-ci}
\end{table}

\begin{pendienteenv}
resultados medibles (tiempos, cobertura, volumen) a medida que se ejecuten las
pruebas pendientes de la §4.
\end{pendienteenv}

\subsection{Limitación operativa: capa gratuita de NVIDIA Build}

Las pruebas que involucran el enriquecimiento con IA están sujetas a la cuota de la
capa gratuita de \textbf{NVIDIA Build} (modelos en la nube, 40 peticiones por minuto
global): no se puede planificar una batería intensiva de pruebas sobre el pipeline
completo sin agotar la cuota ni saturar la latencia. Por ello, las corridas de prueba
del pipeline deben planificarse por muestras (subconjuntos de fuentes o de fichas) y
aprovechar el mecanismo de checkpoints, que evita reprocesar lo ya completado. Si la
ejecución se interrumpe por la cuota, se reanuda sin repetir trabajo.

\begin{pendienteenv}
calendario de corridas de prueba (fechas y tamaño de muestra) una vez que el equipo
defina la ventana de pruebas del bloque 1.
\end{pendienteenv}

\end{document}
